\paragraph{(U) 型\trBrane{}与全局单值。}\label{sub:type-u-branes-global-monodromy}
照前文，设 $\boldsymbol{\gamma}$ 是\trVanishingPath{}的一个\trDistinguishedBasis。给每个 $\gamma_k$ 配上 $\widetilde{E}$ 中与之相伴的 (B) 型\trEquivariant\trBrane{}$\widetilde{\Delta}^{\#}_k$，其中\trGrading{}的选取是任意的。此外，设 $L^{\#}$ 是 $E$ 中一个含于 $E^{\mathrm{shrink}}$ 之内的\trExactLagrangianBrane，并令 $\widetilde{L}^{\#}=\widetilde{L}^{\#}_{+}\cup\widetilde{L}^{\#}_{-}$ 为与之对应的 (U) 型\trEquivariant\trBrane。
\begin{lemma}\label{thm:type-u-shift-relation}
在\trDerivedCategory{}$D(\widetilde{\scr{A}})=H^{0}(Tw\widetilde{\scr{A}})$ 中，
\begin{equation}\label{eq:type-u-twist-shift}
T_{\widetilde{\Delta}^{\#}_1}\cdots T_{\widetilde{\Delta}^{\#}_m}
\bigl(\widetilde{L}^{\#}_{+}\bigr)\cong S\bigl(\widetilde{L}^{\#}_{-}\bigr).
\end{equation}
\end{lemma}
\begin{proof}
略去\trBraneStructure{}与\trSuspension{}$S$ 的一个简化版本可以这样来解释。由\cref{thm:twist-equals-tau}，\eqref{eq:type-u-twist-shift}左边的\trObject{}\trIsomorphic{}于由\trLagrange\trSubmanifold{}$\tau_{\widetilde{\Delta}_1}\cdots\tau_{\widetilde{\Delta}_m}(\widetilde{L}_{+})$ 所给出的那个\trObject。\cref{thm:matching-vanishing-isotopy}告诉我们：（在\trFrame\trIsotopy{}的意义下）诸 $\widetilde{\Delta}_k$ 正是把 $\gamma_k$ 看成 $\Psi\colon\widetilde{\scr{E}}\to\scr{W}$ 的一条\trVanishingPath{}时与之相伴的那些\trVanishingCycle，从而它们的\trDehnTwist{}都是 $\Psi$ 的\trMonodromy。由此可知，全部 $\tau_{\widetilde{\Delta}_k}$ 的\trComposition{}在 $\operatorname{Aut}(\widetilde{E})$ 中\trIsotopy{}于 $\Psi$ 沿 $\scr{W}$ 中半径为 $1-3\delta$ 的\trCircle{}一周的\trMonodromy。由\cref{thm:monodromy-involution-agreement}，那个\trMonodromy{}在子集 $\widetilde{E}^{\mathrm{shrink}}$ 上的限制等于 $\iota_{\widetilde{E}}$（按定义，该子集包含诸 $\widetilde{L}_{\pm}$），因而把 $\widetilde{L}_{+}$ 映到 $\widetilde{L}_{-}$。
为了得到完整的陈述，可以更加细心一点地论证如下。$T_{\widetilde{\Delta}^{\#}_k}$ 在\trObject{}上的作用对应于配备了\trCanonical\trGrading{}的 $\tau_{\widetilde{\Delta}_k}$ 的作用，而由\trPicard--\trLefschetz{}公式的\trGraded{}版本，后者也就是与 $\Psi$ 的相应\trMonodromy{}自然相伴的那个\trGrading。由\cref{thm:monodromy-natural-grading}，$\widetilde{L}^{\#}_{+}=(\widetilde{L}_{+},\tilde{\alpha}^{\#}_{+},P^{\#}_{+})$ 在这个\trMonodromy{}之下的像是 $(\iota_{\widetilde{E}}(\widetilde{L}_{+}),\iota_{\widetilde{E}}^{*}\tilde{\alpha}^{\#}_{+}-1,\iota_{\widetilde{E}}^{*}P^{\#}_{+})$。它与 $\widetilde{L}^{\#}_{-}$ 相差\trGrading{}中的一个常值 $-1$，以及\trPinStructure{}中的一个 $\otimes\lambda^{\mathrm{top}}(T\widetilde{L}_{-})$\footnote{译注：原书此处印作 $\lambda^{\mathrm{top}}(TL_{-})$；疑应为 $\lambda^{\mathrm{top}}(T\widetilde{L}_{-})$（本处所论\trBrane{}是 $\widetilde{E}$ 中的 $\widetilde{L}^{\#}_{-}$，而\trSuspension{}操作 \eqref{eq:shift-operation} 所含的那个因子正取自该\trBrane{}底\trLagrange\trSubmanifold{}的\trTangentBundle，$E$ 中的 $L$ 并无 $\pm$ 分量）；此处径改。——译者注}。但这个差别恰好就是\cref{sub:lagrangian-branes}中所定义的\trSuspension{}操作在\trLagrange\trBrane{}上的作用；我们在\cref{sub:shift-and-a-infinity-structure}中检验过，它确实\trQuasiRepresent{}抽象的（代数的）\trSuspension。
\end{proof}
\begin{remark}\label{thm:shift-restores-orientation}
如果把 $S$ 看成一个反转\trLagrange\trSubmanifold{}\trOrientation{}的操作，那么它在 \eqref{eq:type-u-twist-shift} 中的出现可能显得令人困惑。事实上，相反的\trOrientation{}已经内置在 $\widetilde{L}^{\#}_{\pm}$ 的定义之中（见\cref{thm:square-roots-natural-orientation}），而\trSuspension{}只不过是恢复一个共同的选择罢了。
\end{remark}
\begin{proposition}\label{thm:type-u-in-generated-subcategory}
$\widetilde{L}^{\#}$ 落在 $Tw\scr{F}(\pi)$ 中由 $\widetilde{\Delta}^{\#}_1,\ldots,\widetilde{\Delta}^{\#}_m$ 生成的那个\trTriangulated\trSubcategory{}之内。
\end{proposition}
\begin{proof}
这是\cref{thm:equivariant-split-generation}的一个直截了当的应用。$\widetilde{L}^{\#}$ 是一个\trEquivariant\trBrane，因而是 $\widetilde{\scr{A}}$ 中在\trInvolution{}$\epsilon_{\widetilde{\scr{A}}}$ 之下不动的\trObject。按定义，它\trIsomorphic{}于 $\widetilde{L}^{\#}_{+}$ 与 $\widetilde{L}^{\#}_{-}=\epsilon_{\widetilde{\scr{A}}}(\widetilde{L}^{\#}_{+})$ 的\trDirectSum。此外，$HF^{*}(\widetilde{L}^{\#}_{+},\widetilde{L}^{\#}_{-})=0$，因为底下的\trLagrange\trSubmanifold{}互不相交。诸 $\widetilde{\Delta}^{\#}_k$ 同样是 $\epsilon_{\widetilde{\scr{A}}}$-不动的\trObject，并且由 \eqref{eq:cf-delta-vanishing-iso}，对一切 $k<l$，$HF^{*}(\widetilde{\Delta}^{\#}_k,\widetilde{\Delta}^{\#}_l)$ 都是 $\ZZ/2$-不变的。\cref{thm:type-u-shift-relation}恰好就是所需要的关系 \eqref{eq:iterated-twist-split}，而且按定义 $\widetilde{\scr{A}}^{\mathrm{inv}}=\scr{F}(\pi)$。
\end{proof}
